% Options for packages loaded elsewhere
% Options for packages loaded elsewhere
\PassOptionsToPackage{unicode}{hyperref}
\PassOptionsToPackage{hyphens}{url}
\PassOptionsToPackage{dvipsnames,svgnames,x11names}{xcolor}
%
\documentclass[
  12pt,
]{article}
\usepackage{xcolor}
\usepackage[top=2cm,bottom=2cm,left=2.5cm,right=2.5cm]{geometry}
\usepackage{amsmath,amssymb}
\setcounter{secnumdepth}{5}
\usepackage{iftex}
\ifPDFTeX
  \usepackage[T1]{fontenc}
  \usepackage[utf8]{inputenc}
  \usepackage{textcomp} % provide euro and other symbols
\else % if luatex or xetex
  \usepackage{unicode-math} % this also loads fontspec
  \defaultfontfeatures{Scale=MatchLowercase}
  \defaultfontfeatures[\rmfamily]{Ligatures=TeX,Scale=1}
\fi
\usepackage{lmodern}
\ifPDFTeX\else
  % xetex/luatex font selection
\fi
% Use upquote if available, for straight quotes in verbatim environments
\IfFileExists{upquote.sty}{\usepackage{upquote}}{}
\IfFileExists{microtype.sty}{% use microtype if available
  \usepackage[]{microtype}
  \UseMicrotypeSet[protrusion]{basicmath} % disable protrusion for tt fonts
}{}
\makeatletter
\@ifundefined{KOMAClassName}{% if non-KOMA class
  \IfFileExists{parskip.sty}{%
    \usepackage{parskip}
  }{% else
    \setlength{\parindent}{0pt}
    \setlength{\parskip}{6pt plus 2pt minus 1pt}}
}{% if KOMA class
  \KOMAoptions{parskip=half}}
\makeatother
% Make \paragraph and \subparagraph free-standing
\makeatletter
\ifx\paragraph\undefined\else
  \let\oldparagraph\paragraph
  \renewcommand{\paragraph}{
    \@ifstar
      \xxxParagraphStar
      \xxxParagraphNoStar
  }
  \newcommand{\xxxParagraphStar}[1]{\oldparagraph*{#1}\mbox{}}
  \newcommand{\xxxParagraphNoStar}[1]{\oldparagraph{#1}\mbox{}}
\fi
\ifx\subparagraph\undefined\else
  \let\oldsubparagraph\subparagraph
  \renewcommand{\subparagraph}{
    \@ifstar
      \xxxSubParagraphStar
      \xxxSubParagraphNoStar
  }
  \newcommand{\xxxSubParagraphStar}[1]{\oldsubparagraph*{#1}\mbox{}}
  \newcommand{\xxxSubParagraphNoStar}[1]{\oldsubparagraph{#1}\mbox{}}
\fi
\makeatother


\usepackage{longtable,booktabs,array}
\usepackage{calc} % for calculating minipage widths
% Correct order of tables after \paragraph or \subparagraph
\usepackage{etoolbox}
\makeatletter
\patchcmd\longtable{\par}{\if@noskipsec\mbox{}\fi\par}{}{}
\makeatother
% Allow footnotes in longtable head/foot
\IfFileExists{footnotehyper.sty}{\usepackage{footnotehyper}}{\usepackage{footnote}}
\makesavenoteenv{longtable}
\usepackage{graphicx}
\makeatletter
\newsavebox\pandoc@box
\newcommand*\pandocbounded[1]{% scales image to fit in text height/width
  \sbox\pandoc@box{#1}%
  \Gscale@div\@tempa{\textheight}{\dimexpr\ht\pandoc@box+\dp\pandoc@box\relax}%
  \Gscale@div\@tempb{\linewidth}{\wd\pandoc@box}%
  \ifdim\@tempb\p@<\@tempa\p@\let\@tempa\@tempb\fi% select the smaller of both
  \ifdim\@tempa\p@<\p@\scalebox{\@tempa}{\usebox\pandoc@box}%
  \else\usebox{\pandoc@box}%
  \fi%
}
% Set default figure placement to htbp
\def\fps@figure{htbp}
\makeatother


% definitions for citeproc citations
\NewDocumentCommand\citeproctext{}{}
\NewDocumentCommand\citeproc{mm}{%
  \begingroup\def\citeproctext{#2}\cite{#1}\endgroup}
\makeatletter
 % allow citations to break across lines
 \let\@cite@ofmt\@firstofone
 % avoid brackets around text for \cite:
 \def\@biblabel#1{}
 \def\@cite#1#2{{#1\if@tempswa , #2\fi}}
\makeatother
\newlength{\cslhangindent}
\setlength{\cslhangindent}{1.5em}
\newlength{\csllabelwidth}
\setlength{\csllabelwidth}{3em}
\newenvironment{CSLReferences}[2] % #1 hanging-indent, #2 entry-spacing
 {\begin{list}{}{%
  \setlength{\itemindent}{0pt}
  \setlength{\leftmargin}{0pt}
  \setlength{\parsep}{0pt}
  % turn on hanging indent if param 1 is 1
  \ifodd #1
   \setlength{\leftmargin}{\cslhangindent}
   \setlength{\itemindent}{-1\cslhangindent}
  \fi
  % set entry spacing
  \setlength{\itemsep}{#2\baselineskip}}}
 {\end{list}}
\usepackage{calc}
\newcommand{\CSLBlock}[1]{\hfill\break\parbox[t]{\linewidth}{\strut\ignorespaces#1\strut}}
\newcommand{\CSLLeftMargin}[1]{\parbox[t]{\csllabelwidth}{\strut#1\strut}}
\newcommand{\CSLRightInline}[1]{\parbox[t]{\linewidth - \csllabelwidth}{\strut#1\strut}}
\newcommand{\CSLIndent}[1]{\hspace{\cslhangindent}#1}



\setlength{\emergencystretch}{3em} % prevent overfull lines

\providecommand{\tightlist}{%
  \setlength{\itemsep}{0pt}\setlength{\parskip}{0pt}}



 


\usepackage{array}
\makeatletter
\@ifpackageloaded{tcolorbox}{}{\usepackage[skins,breakable]{tcolorbox}}
\@ifpackageloaded{fontawesome5}{}{\usepackage{fontawesome5}}
\definecolor{quarto-callout-color}{HTML}{909090}
\definecolor{quarto-callout-note-color}{HTML}{0758E5}
\definecolor{quarto-callout-important-color}{HTML}{CC1914}
\definecolor{quarto-callout-warning-color}{HTML}{EB9113}
\definecolor{quarto-callout-tip-color}{HTML}{00A047}
\definecolor{quarto-callout-caution-color}{HTML}{FC5300}
\definecolor{quarto-callout-color-frame}{HTML}{acacac}
\definecolor{quarto-callout-note-color-frame}{HTML}{4582ec}
\definecolor{quarto-callout-important-color-frame}{HTML}{d9534f}
\definecolor{quarto-callout-warning-color-frame}{HTML}{f0ad4e}
\definecolor{quarto-callout-tip-color-frame}{HTML}{02b875}
\definecolor{quarto-callout-caution-color-frame}{HTML}{fd7e14}
\makeatother
\makeatletter
\@ifpackageloaded{caption}{}{\usepackage{caption}}
\AtBeginDocument{%
\ifdefined\contentsname
  \renewcommand*\contentsname{Table of contents}
\else
  \newcommand\contentsname{Table of contents}
\fi
\ifdefined\listfigurename
  \renewcommand*\listfigurename{List of Figures}
\else
  \newcommand\listfigurename{List of Figures}
\fi
\ifdefined\listtablename
  \renewcommand*\listtablename{List of Tables}
\else
  \newcommand\listtablename{List of Tables}
\fi
\ifdefined\figurename
  \renewcommand*\figurename{Figure}
\else
  \newcommand\figurename{Figure}
\fi
\ifdefined\tablename
  \renewcommand*\tablename{Table}
\else
  \newcommand\tablename{Table}
\fi
}
\@ifpackageloaded{float}{}{\usepackage{float}}
\floatstyle{ruled}
\@ifundefined{c@chapter}{\newfloat{codelisting}{h}{lop}}{\newfloat{codelisting}{h}{lop}[chapter]}
\floatname{codelisting}{Listing}
\newcommand*\listoflistings{\listof{codelisting}{List of Listings}}
\makeatother
\makeatletter
\makeatother
\makeatletter
\@ifpackageloaded{caption}{}{\usepackage{caption}}
\@ifpackageloaded{subcaption}{}{\usepackage{subcaption}}
\makeatother
\usepackage{bookmark}
\IfFileExists{xurl.sty}{\usepackage{xurl}}{} % add URL line breaks if available
\urlstyle{same}
\hypersetup{
  pdftitle={Parallel Computing Project Report},
  pdfauthor={Leonardo Bertani; Rafael Nichel Pearl},
  colorlinks=true,
  linkcolor={blue},
  filecolor={Maroon},
  citecolor={Blue},
  urlcolor={Blue},
  pdfcreator={LaTeX via pandoc}}


\title{Parallel Computing Project Report}
\usepackage{etoolbox}
\makeatletter
\providecommand{\subtitle}[1]{% add subtitle to \maketitle
  \apptocmd{\@title}{\par {\large #1 \par}}{}{}
}
\makeatother
\subtitle{Matrix Multiplication}
\author{Leonardo Bertani \and Rafael Nichel Pearl}
\date{}
\begin{document}
\maketitle


\begin{center}
  Academic Year 2025/2026
  \vspace{1em}
\end{center}

\section{Introduction}\label{introduction}

Matrix multiplication is one of the most fundamental operations in
linear algebra, used for a wide range of applications, including
computer graphics, machine learning, and high-performance scientific
simulations.

Given the importance of that operation, several hardware manufacturers
have developed specialized co-processors and embedded calculators to
execute these matrix routines directly at the hardware level.

However, this report focuses on the software-level implementation and
optimization of this workload. We first explore a baseline sequential
approach, and then we will try to parallelize that algorithm using MPI.

\begin{tcolorbox}[enhanced jigsaw, rightrule=.15mm, leftrule=.75mm, colback=white, toprule=.15mm, breakable, arc=.35mm, bottomrule=.15mm, left=2mm, opacityback=0, colframe=quarto-callout-note-color-frame]

Although the general procedure for matrix multiplication is valid for
any pair of matrices where the number of columns in matrix \(A\) is
equal to the number of rows in matrix \(B\), this report focuses on
square matrices of size \(N \times N\). This choice is motivated by
structural and algebraic symmetry, which simplifies the theoretical
analysis and allows for an optimal, balanced mapping of the
computational workload onto the parallel processor topology.

\end{tcolorbox}

\section{Sequential algorithm}\label{sequential-algorithm}

\subsection{Problem definition}\label{problem-definition}

Given two square matrices \(A\) and \(B\) of dimension \(N \times N\),
the goal is to compute their product \(C = A \times B\), where \(C\) is
also an \(N \times N\) matrix.

\begin{table}[h]
\centering
\begin{tabular}{m{0.68\textwidth} m{0.28\textwidth}}
The element $C_{ij}$ of the resulting matrix is computed by taking the dot product of the $i$-th row of $A$ and the $j$-th column of $B$. & 
$C_{ij} = \sum_{k=0}^{N-1} A_{ik} \cdot B_{kj}$
\end{tabular}
\end{table}

\subsection{Algorithm and complexity}\label{algorithm-and-complexity}

In the following figures, we present the sequential algorithm for matrix
multiplication in pseudocode and C code. This follows the approach
described in the previous section step-by-step.

\begin{center}
\begin{minipage}[t]{0.46\textwidth}
\subsubsection*{Pseudocode}
\begin{small}
\begin{verbatim}
Sequential MM(A, B, C, N):
  Initialize C with all zeros
  For i = 0 To N - 1:
    For j = 0 To N - 1:
      For k = 0 To N - 1:
        C[i][j] += A[i][k] * B[k][j]
\end{verbatim}
\end{small}
\end{minipage}
\hfill
\begin{minipage}[t]{0.48\textwidth}
\subsubsection*{C Code}
\begin{small}
\begin{verbatim}
void matrix_multiply(double *A, 
        double *B, double *C, int N) {
    for (int i = 0; i < N; i++) {
        for (int j = 0; j < N; j++) {
            C[i][j] = 0.0;
            for (int k = 0; k < N; k++) {
                C[i][j] += A[i][k] * B[k][j];
}}}}
\end{verbatim}
\end{small}
\end{minipage}
\end{center}

\vspace{0.01\textheight}

As we can see in the code above, the algorithm is composed of three
nested loops. The outer two loops iterate over the rows and columns of
the resulting matrix \(C\). The inner loop iterates through the elements
of the \(i\)-th row of matrix \(A\) and the \(j\)-th column of matrix
\(B\) simultaneously to compute their dot product, following the
standard row-by-column multiplication strategy.

About the time complexity, this algorithm has a time complexity of
\(\Theta(N^3)\). Then, regarding the space complexity, the algorithm
requires \(\Theta(N^2)\) auxiliary space to store the input matrices
\(A\), \(B\), and the resulting matrix \(C\) in memory.

\section{Parallel design}\label{parallel-design}

Some considerations must be addressed before trying to parallelize the
algorithm. In particular, we can analyze the dependency of the algorithm
and the communication pattern, to ensure an optimal parallel
implementation.

\subsection{Dependency analysis}\label{dependency-analysis}

Before parallelizing the algorithm, it is critical to analyze how data
flows through the hardware. While mathematically straightforward, the
sequential \(i\text{-}j\text{-}k\) loop order introduces a significant
performance bottleneck due to its asymmetric memory access pattern:

\begin{itemize}
\tightlist
\item
  Matrix \(A\) is accessed row-by-row (\(A[i][k]\)), which perfectly
  aligns with the row-major storage of C arrays. This benefits from
  excellent \emph{spatial locality} in the cache.
\item
  Matrix \(B\) is traversed column-by-column (\(B[k][j]\)). This creates
  a stride-\(N\) access pattern, loading an entire cache line but only
  using a single element before moving to the next row. As \(N\) grows,
  this triggers a massive rate of \emph{cache misses}.
\end{itemize}

A dependency analysis reveals a high degree of data parallelism when
mapping this workload to multiple execution units. The iterations of the
two outer loops (\texttt{i} and \texttt{j}) are completely independent,
since computing any element \(C_{ij}\) does not depend on, nor update,
any other element of the output matrix \(C\).

From a practical implementation standpoint, this independent structure
allows for a straightforward \emph{1D Domain Decomposition}. The
algorithm can be parallelized by assigning a contiguous subset of rows
of matrix \(A\) to each thread or process. Each execution unit can then
compute its corresponding rows of matrix \(C\) independently.

\subsection{Communication pattern}\label{communication-pattern}

While the 1D decomposition eliminates the need for synchronization in a
shared-memory system (e.g., OpenMP) where all threads access a single
global RAM, it introduces severe scalability limits in a
distributed-memory system (e.g., MPI). In a distributed environment, a
1D approach requires a communication pattern where every single process
must replicate and hold the entire matrix \(B\) in its private memory,
saturating network bandwidth via heavy broadcast operations.

To overcome this network bottleneck, the workload can be restructured
into a \emph{2D Domain Decomposition}, creating a logical 2D grid of
processes. Under this scheme, each process is responsible for computing
only a localized submatrix (block) of \(C\), requiring access to just a
submatrix of \(A\) and a submatrix of \(B\).

However, in a naive 2D parallel implementation, a processor still needs
to perform collective communication (broadcasts) across its entire
virtual row and column to fetch the required blocks of \(A\) and \(B\)
over time. This communication overhead heavily bounds scalability,
driving the need for a fully localized, communication-minimizing
topology configuration, such as Cannon's Algorithm.

\section{Topology-aware
implementation}\label{topology-aware-implementation}

\subsection{Topologies considered}\label{topologies-considered}

In a matrix multiplication algorithm, the computation naturally operates
on a fixed geometric space defined by its nested loops. When mapping
these loop structures to parallel hardware, closed topologies like the
1D Ring and 2D Torus offer significant code-level and algorithmic
advantages over open topologies like the Array or Mesh.

For instance, closed topologies are perfectly symmetric, meaning every
node has the exact same number of neighbors. This structural uniformity
allows every processing unit to execute the identical chunk of code
without requiring complex conditional statements to handle boundary
exceptions. Additionally, physically bridging the boundaries cuts the
maximum routing length in half, allowing the system to optimize the
direction in which data is sent and drastically reducing latency during
the data-shifting phases.

\subsection{Embedding the communication
pattern}\label{embedding-the-communication-pattern}

When embedding the logical communication pattern of the matrix
multiplication loops onto the hardware, the choice between a 1D Ring and
a 2D Torus fundamentally alters how data flows through the network and
how the code is structured.

\subsubsection{1D Ring Embedding}\label{d-ring-embedding}

In a 1D Ring configuration, the multi-dimensional problem must be
compressed into a single linear dimension. The matrix is typically
divided into horizontal slices (row blocks), and each node communicates
only with its two linear neighbors.

While this topology naturally supports the sequential rotation of one
matrix axis, it creates a severe structural bottleneck for the remaining
dimensions. Because the hardware lacks a second physical dimension, any
data reuse or shifting required by the nested loops along the column
axis cannot be localized. Data must either be globally replicated across
the entire ring or bounced across multiple intermediate nodes, scaling
the communication latency linearly with the number of processors.

\subsubsection{2D Torus Embedding}\label{d-torus-embedding}

The 2D Torus provides a perfect geometric match for the nested loop
structure of matrix multiplication, as it maps the execution directly
onto a two-dimensional physical grid of processors.

By having four physical connections per node with wrap-around links (so
the degree is 4 for every node, also on the boundaries), the
communication pattern achieves optimal spatial locality in both
dimensions simultaneously. The cyclic data shifts required by the outer
and inner loops are executed through independent, concurrent paths. Data
moves exactly one hop to an immediate neighbor in a single step,
preventing network contention. This multi-directional routing eliminates
the non-local data traffic inherent to the 1D Ring, keeping the
communication overhead constant and ensuring maximum computational
efficiency as the processor grid scales.

\section{MPI implementation}\label{mpi-implementation}

\subsection{Language and build}\label{language-and-build}

For this project, the C programming language and the MPI library were
selected to implement the parallel matrix multiplication algorithm.

Initially, Python was considered as an alternative. However, standard
Python is inherently slow for low-level memory operations, while
high-level numerical libraries like NumPy abstract away and
automatically optimize matrix computations under the hood. Since NumPy
handles memory allocation and parallelism internally, it prevents the
explicit control and manual implementation of low-level data movement
and communication patterns required by this project. Therefore, C was
chosen to ensure maximum execution speed, precise memory management, and
total control over the MPI communication primitives.

\subsection{Topology realisation in
MPI}\label{topology-realisation-in-mpi}

To implement a topology-aware algorithm like Cannon's, the abstract
concept of a 2D Torus must be explicitly declared within the MPI
execution environment.

By default, MPI initializes processes in a flat, linear communicator
(MPI\_COMM\_WORLD), which lacks any geometric awareness.

The concrete realization of the 2D Torus topology is achieved through
three core functional phases in the code: 1. Grid Creation and
Periodicity (\texttt{MPI\_Cart\_create}): 1D \(\rightarrow\) 2D of size
\(\sqrt{P} \times \sqrt{P}\) plus wrap around edges using the
periodicity flags (periods = \{1, 1\}) 2. Automated Neighbor Discovery
(\texttt{MPI\_Cart\_shift}): to automate the identification of
communication partners, returning the ranks of the neighbors in the
specified direction 3. Deadlock-Free Data Swapping
(\texttt{MPI\_Sendrecv}): combining the outgoing send and incoming
receive operations into a single, atomic transaction.

\section{Experimental setup}\label{experimental-setup}

\section{Results}\label{results}

\section{Discussion}\label{discussion}

\subsection{Limitations}\label{limitations}

This algorithm strictly requires that the global matrices \(A\) and
\(B\) are square (\(N \times N\)). Furthermore, for the parallel
implementation of Cannon's algorithm, the total number of MPI processes
(\(P\)) must be a perfect square, meaning that the logical grid size
\(q = \sqrt{P}\) must be an integer.

A major practical constraint is that the global matrix dimension \(N\)
must be perfectly divisible by \(q\) (\(N \pmod q = 0\)).

If these structural conditions are not met, the workload cannot be
evenly distributed across the 2D torus topology without implementing
complex padding or masking techniques, which would significantly degrade
parallel efficiency.

Additionally, as a standard 2D matrix multiplication method, Cannon's
algorithm is bound to a strict communication lower bound of
\(\Omega(N^2 / \sqrt{P})\) words sent or received per processor.

\subsection{Possible improvements}\label{possible-improvements}

We could explore alternative parallelization strategies to reach a more
flexible solution in handling non-square matrices or processor grids
that are not perfectly square.

Furthermore, communication overhead can be mitigated by upgrading from a
2D topology to a 3D parallel matrix multiplication algorithm. While 2D
algorithms like Cannon's are strictly bound by a communication lower
bound of \(\Omega(N^2 / \sqrt{P})\) words sent or received per
processor, 3D approaches replicate matrix data across a
three-dimensional processor grid. This technique reduces the total
network communication volume per processor to \(\Omega(N^2 / P^{2/3})\)
, though it does so at the expense of an increased local memory
footprint. {[}1{]}

However, the asymptotic reduction in communication achieved by a 3D
layout compared to a 2D algorithm like Cannon is limited to a factor of
only \(P^{1/6}\), meaning that the overhead of data replication may not
always be profitable for standard cluster sizes. {[}1{]}

To bridge this gap, a 2.5D variant can be introduced {[}2{]}. This
method incorporates a tunable replication factor that dynamically scales
according to the available RAM to reach optimal, communication-avoiding
bounds. Additionally, adapting these advanced mapping strategies to
higher-dimensional physical interconnects, such as hypercube network
topologies~{[}3{]}, can further minimize routing latency and maximize
bandwidth, yielding massive speedups on thousands of processing cores.

\section{Conclusions}\label{conclusions}

\section*{References}\label{references}
\addcontentsline{toc}{section}{References}

\phantomsection\label{refs}
\begin{CSLReferences}{0}{0}
\bibitem[\citeproctext]{ref-irony2004}
\CSLLeftMargin{{[}1{]} }%
\CSLRightInline{D. Irony, S. Toledo, and A. Tiskin, {``Communication
lower bounds for distributed-memory matrix multiplication,''}
\emph{Journal of Parallel and Distributed Computing}, vol. 64, no. 9,
pp. 1017--1026, 2004, doi:
\href{https://doi.org/10.1016/j.jpdc.2004.03.021}{10.1016/j.jpdc.2004.03.021}.}

\bibitem[\citeproctext]{ref-solomonik2011}
\CSLLeftMargin{{[}2{]} }%
\CSLRightInline{E. Solomonik and J. Demmel, {``Communication-optimal
parallel 2.5D matrix multiplication and LU factorization algorithms,''}
in \emph{Euro-par 2011 parallel processing}, in Lecture notes in
computer science, vol. 6853. Berlin, Heidelberg: Springer, 2011, pp.
90--109. doi:
\href{https://doi.org/10.1007/978-3-642-23397-5_10}{10.1007/978-3-642-23397-5\_10}.}

\bibitem[\citeproctext]{ref-berntsen1989}
\CSLLeftMargin{{[}3{]} }%
\CSLRightInline{J. Berntsen, {``Communication efficient matrix
multiplication on hypercubes,''} \emph{Parallel Computing}, vol. 12, no.
3, pp. 335--342, 1989, doi:
\href{https://doi.org/10.1016/0167-8191(89)90031-8}{10.1016/0167-8191(89)90031-8}.}

\bibitem[\citeproctext]{ref-wiki:matrix_multiplication}
\CSLLeftMargin{{[}4{]} }%
\CSLRightInline{Wikipedia contributors, {``Matrix multiplication.''}
\url{https://en.wikipedia.org/wiki/Matrix_multiplication}, 2026.}

\end{CSLReferences}




\end{document}
